\documentclass[11pt]{article}
\usepackage[margin=1in]{geometry}
\usepackage{amsmath, amsthm, amssymb}
\usepackage{hyperref}
\usepackage{enumitem}
\usepackage{microtype}
\sloppy

\newtheorem*{theorem*}{Theorem}
\newtheorem*{keylemma*}{Key lemma}

\title{The Iwasawa Decomposition of $GL_n(\mathbb{R})$}
\author{Jiho Lee \\ Math 157}
\date{April 28, 2026}

\begin{document}
\maketitle

\section*{Motivation}

The Iwasawa decomposition expresses every invertible real $n \times n$
matrix as a product of three matrices, each drawn from a distinguished
subgroup. Lang treats it in Appendix~II of \emph{Linear Algebra} (3rd ed.,
1987), titled ``Iwasawa Decomposition and Others.'' The appendix's stated
purpose is to introduce the language of group theory through the running
example of $SL_n$ (the group of $n \times n$ real matrices with
determinant~$1$), which had not been used in the rest of the book; Lang
chooses the Iwasawa decomposition as the vehicle, describing it
explicitly as a way to ``express Theorem~2.1 of Chapter~V in the context
of groups and subgroups.'' That theorem is the Gram--Schmidt
orthonormalization theorem, so the Iwasawa decomposition is, for Lang,
the group-theoretic restatement of Gram--Schmidt: the assertion that the
product map $U \times A \times K \to SL_n$ is a bijection, where $U$ is
the upper-unipotent subgroup, $A$ the positive diagonal subgroup, and
$K$ the orthogonal subgroup. From there the appendix moves on to the
broader content suggested by ``Others'': the actions of
$SL_2(\mathbb{R})$ and $SL_2(\mathbb{C})$ on the hyperbolic upper
half-plane and upper half-space, the conjugation action with the first
rudiments of Lie algebras, and the complementary polar decomposition
$G = KAK$. The Iwasawa decomposition itself is the matrix face of a
more general phenomenon in Lie theory: any real semisimple Lie group
$G$ admits a diffeomorphism $G \cong K \times A \times N$, where $K$
is a maximal compact subgroup, $A$ a maximal split abelian subgroup,
and $N$ a maximal unipotent subgroup. The formalization here works
with the slightly larger group $GL_n(\mathbb{R})$ (no determinant
constraint); the proof carries over from Lang's argument with no
essential change.

\section*{Statement and proof}

Write $M_n$ for the type of $n \times n$ real matrices. A matrix is
\emph{positive diagonal} if it is diagonal with strictly positive
entries, and \emph{upper unipotent} if it is upper triangular with ones
on the diagonal.

\begin{theorem*}[Iwasawa decomposition for $GL_n(\mathbb{R})$]
For every $g \in M_n$ with $\det g \neq 0$ there exist unique
$k, a, u \in M_n$ such that $g = k \cdot a \cdot u$, where $k$ is
orthogonal, $a$ is positive diagonal, and $u$ is upper unipotent. (Lang
states the bijection in the form $g = uak$; the two orderings are
equivalent by inversion since each of $U$, $A$, $K$ is closed under
inverses.)
\end{theorem*}

\subsection*{Existence via Gram--Schmidt}

Following Lang's Appendix~II, view column $i$ of $g$ as
$g^{(i)} \in \mathbb{R}^n$ with $(g^{(i)})_j = g_{ji}$. Since
$\det g \neq 0$, the columns are linearly independent. Apply
Gram--Schmidt: produce vectors $\tilde{e}_i$, the components of
$g^{(i)}$ orthogonal to the span of the earlier columns, and let
$e_i = \tilde{e}_i / \|\tilde{e}_i\|$. Form $Q$ with $i$-th column
$e_i$; orthonormality of $(e_i)$ is exactly $Q^T Q = I$, so $Q$ is
orthogonal. Set $R = Q^T g$, equivalently
$R_{ij} = \langle e_i, g^{(j)} \rangle$. Two facts about $R$ follow
from Gram--Schmidt: (i)~for $j < i$, the $i$-th Gram--Schmidt vector
is orthogonal to $g^{(j)}$ (which lies in the span of earlier
columns), so $R_{ij} = 0$ and $R$ is upper triangular;
(ii)~expanding $g^{(i)}$ as $\tilde{e}_i$ plus projections onto the
earlier $\tilde{e}_k$ and pairing with $\tilde{e}_i$ collapses to
$\langle \tilde{e}_i, g^{(i)} \rangle = \|\tilde{e}_i\|^2$, hence
$R_{ii} = \|\tilde{e}_i\| > 0$. Now split $R = D \cdot U$ with $D$
the diagonal of $R$ and $U = D^{-1} R$. Then $D$ is positive
diagonal, and $U$ is upper unipotent
($U_{ii} = R_{ii}^{-1} \cdot R_{ii} = 1$). Combining, $g = QR = Q(DU)$,
so $k = Q$, $a = D$, $u = U$ is the required factorization.

\subsection*{Uniqueness via a key lemma}

Uniqueness rests on:

\begin{keylemma*}
If $M$ is orthogonal, upper triangular, and has positive diagonal, then
$M = I$.
\end{keylemma*}

\begin{proof}
Orthogonality gives $M^{-1} = M^T$. The inverse of an upper-triangular
invertible matrix is upper triangular, so $M^T$ is upper triangular,
equivalently $M$ is lower triangular. With both upper and lower
triangular, $M$ is diagonal. From $M M^T = I$ each $M_{ii}^2 = 1$, and
positivity selects $M_{ii} = 1$.
\end{proof}

To finish uniqueness, suppose $g = k_1 a_1 u_1 = k_2 a_2 u_2$ and set
$M = k_2^T k_1$. Then $M$ is orthogonal (product of orthogonals), and
rearranging gives $M = a_2 u_2 u_1^{-1} a_1^{-1}$, which is upper
triangular with positive diagonal (the diagonal of a product of
upper-triangular matrices is the product of diagonals; positivity
follows from positivity of $a_1, a_2$ and unipotence of $u_1, u_2$).
The lemma forces $M = I$, so $k_1 = k_2$. Cancelling gives
$a_1 u_1 = a_2 u_2$, hence $a_1 = a_2(u_2 u_1^{-1})$; comparing
diagonals forces $a_1 = a_2$ and the unipotent factor to be the
identity, so $u_1 = u_2$.

\section*{The formalization landscape in Mathlib}

Mathlib provides each structural ingredient required by the proof but
does not package the decomposition itself.
\texttt{Mathlib.LinearAlgebra.UnitaryGroup} defines the orthogonal
group on real matrices;
\texttt{Mathlib.Analysis.InnerProductSpace.GramSchmidtOrtho} provides
\texttt{gramSchmidt} and \texttt{gramSchmidtNormed} together with
orthogonality, span, and nonvanishing lemmas; and
\texttt{Mathlib.LinearAlgebra.Matrix.Block} provides the
\texttt{BlockTriangular} predicate together with the result that the
inverse of an upper-triangular invertible matrix is upper triangular.
The closest existing matrix factorization in the library is the LDL
decomposition in \texttt{Mathlib.Analysis.Matrix.LDL}, useful as a
structural template, but neither the Iwasawa decomposition nor the
closely related QR factorization (numerical linear algebra's name for
the same data; see Trefethen--Bau) is packaged anywhere in Mathlib.
The accompanying repository
\url{https://github.com/CuteSurtr/Iwasawa_Decomposition} contains a
complete Lean 4 / Mathlib formalization corresponding to the argument
above; the full formal proof, with explicit references to every
Mathlib lemma used, appears in the repository's \texttt{README.md}.

\section*{References}

\small
\begin{enumerate}[leftmargin=*, labelsep=0.4em, itemsep=0pt, topsep=0pt]
\item S. Lang. \emph{Linear Algebra}, 3rd edition. Springer, 1987.
(Appendix~II, ``Iwasawa Decomposition and Others.'')
\item L. N. Trefethen, D. Bau. \emph{Numerical Linear Algebra}. SIAM,
1997. (Lectures 7--8 on QR.)
\item J. Miao, K. Buzzard, A. Bentkamp.
\texttt{Mathlib.Analysis.InnerProductSpace.GramSchmidtOrtho}.
\item A. Bentkamp. \texttt{Mathlib.Analysis.Matrix.LDL}.
\item The Mathlib community. \url{https://leanprover-community.github.io/}.
\end{enumerate}

\end{document}
